% Transcription — fonds Alexandre Grothendieck, shelfmark 161-3, batch 3 (pages 41-54).
% Established from the facsimile archives/batches/161-3/batch-03.pdf (a mirror of
% https://grothendieck.umontpellier.fr/161-3.pdf), per the skill
% transcribe-grothendieck. All fourteen pages carry mathematics — none skipped.
% Pass: Opus 5 (claude-opus-5), 2026-08-16 — first pass, unchecked against the pages by a human.
\documentclass[11pt,a4paper]{article}
% Le préambule commun à toutes les transcriptions du fonds Grothendieck.
%
% Il tient en une idée : **l'incertitude se compose, elle ne se résout pas.**
% Une transcription qui lisse un mot douteux a détruit la seule chose pour
% laquelle on la faisait — un lecteur doit pouvoir savoir, à chaque endroit,
% ce qui était sur le feuillet et ce qui a été ajouté. Les six macros qui
% suivent sont donc l'appareil critique, et non de la décoration.
%
% Les mêmes commandes sont reconnues par `scripts/render.mjs`, qui produit la
% vue de lecture du site. Une macro ajoutée ici doit l'être aussi là-bas,
% faute de quoi le rendu échouera bruyamment — ce qui est voulu : un rendu
% silencieusement faux est pire qu'un rendu absent.

\ProvidesPackage{grothendieck}[2026/08/08 Transcriptions du fonds Grothendieck]

\usepackage{amsmath,amssymb,amsthm}
\usepackage{mathrsfs}
\usepackage{stmaryrd}
% Les diagrammes commutatifs sont partout dans ce fonds : c'est la langue
% même de la géométrie algébrique de Grothendieck, et une transcription qui
% les rendrait en prose ne transcrirait rien.
\usepackage{tikz-cd}
% Français par défaut — c'est la langue des feuillets — mais l'anglais est
% chargé aussi : la traduction et le résumé sont des documents anglais, et la
% typographie française (espaces avant les deux-points) y serait une faute.
% Ces documents basculent par \selectlanguage{english} en tête de corps.
\usepackage[english,french]{babel}
\usepackage{fontspec}
\usepackage{geometry}
\geometry{a4paper,margin=25mm,marginparwidth=28mm}
\usepackage{marginnote}
\usepackage{xcolor}
% ulem, not soul. `soul`'s \st and \ul are built on a character-by-character
% scan that XeLaTeX's font machinery breaks: the document fails with an
% undefined control sequence rather than degrading. ulem does the same two jobs
% and is Unicode-engine safe. `normalem` stops it hijacking \emph, which the
% transcriptions use for Grothendieck's own emphasis.
\usepackage[normalem]{ulem}
\usepackage{hyperref}
\hypersetup{colorlinks=true,linkcolor=black,urlcolor=black,pdfborder={0 0 0}}

% --- Métadonnées du lot -----------------------------------------------------
% Reprises telles quelles par le script de rendu, qui les affiche en tête de
% la vue de lecture. Elles fixent ce que le fichier prétend couvrir : sans
% elles, une transcription flotte sans point d'attache dans un fonds de seize
% mille feuillets.

\newcommand{\folder}[1]{\gdef\grothFolder{#1}}
\newcommand{\batch}[1]{\gdef\grothBatch{#1}}
\newcommand{\leaves}[2]{\gdef\grothFirst{#1}\gdef\grothLast{#2}}
\newcommand{\foldertitle}[1]{\gdef\grothFoldertitle{#1}}
\gdef\grothFolder{?}\gdef\grothBatch{?}\gdef\grothFirst{?}\gdef\grothLast{?}\gdef\grothFoldertitle{}

% --- Le résumé --------------------------------------------------------------
%
% Il ouvre la lecture modernisée, sous le titre « Résumé », et il est composé
% en petit. Ce n'est pas un ornement : c'est le seul passage du document qui
% ne soit pas des mathématiques, et un lecteur doit pouvoir le reconnaître
% avant de l'avoir lu — puis le sauter s'il connaît déjà le sujet, ou s'y
% arrêter s'il ne le connaît pas. Le filet qui le suit marque la reprise du
% corps.

\newenvironment{resume}
  {\small\setlength{\parskip}{0.45em}}
  {\par\medskip\hrule\medskip}

% --- Mots-clés ---------------------------------------------------------------
%
% En anglais, en clôture du résumé de la lecture modernisée : le vocabulaire
% moderne sous lequel chercher ce que les feuillets construisent. Le manifeste
% du site (`npm run manifest`) les extrait pour étiqueter la cote entière —
% cette ligne est la source unique des tags, il n'y a pas de fichier à côté.
\newcommand{\keywords}[1]{\par\smallskip\noindent
  {\footnotesize\textsc{Keywords} --- \textit{#1}}\par}

% --- L'appareil critique ----------------------------------------------------

% Le numéro de feuillet, en marge. C'est l'ancre qui permet de régler un
% désaccord sur une formule en regardant *un* feuillet plutôt qu'une chemise
% de six cents. Le site s'en sert aussi pour tourner les pages du fac-similé
% quand on fait défiler la transcription.
\newcommand{\leaf}[1]{%
  \marginnote{\footnotesize\color{gray!70}\textsf{#1}}%
  \label{leaf:#1}\ignorespaces}

% Illisible : on ne devine pas. Un mot inventé qui se lit comme les autres est
% le pire résultat possible.
\newcommand{\ill}{\textcolor{red!60!black}{[\,\ldots\,]}}

% Lecture douteuse : le mot est donné, et signalé comme incertain.
\newcommand{\uncertain}[1]{\uline{#1}}

% Ajout de l'éditeur — une lettre manquante, un mot sous-entendu.
\newcommand{\add}[1]{\textcolor{blue!50!black}{[#1]}}

% Biffé par Grothendieck lui-même. Ce qu'il a rayé fait partie du document :
% une rature dit souvent où la pensée a changé de direction.
\newcommand{\struck}[1]{\sout{#1}}

% Note du transcripteur, sur l'état matériel du feuillet.
\newcommand{\note}[1]{\par\smallskip{\small\color{gray!60!black}\textit{[#1]}}\par\smallskip}

% Note marginale de Grothendieck, distincte du corps du texte.
\newcommand{\marginal}[1]{\par\smallskip\noindent\hspace{1em}%
  {\small\color{gray!60!black}#1}\par\smallskip}

% --- Titre ------------------------------------------------------------------

\AtBeginDocument{%
  \begin{center}
    {\large\bfseries Fonds Alexandre Grothendieck --- cote n\textsuperscript{o}~\grothFolder}\\[2pt]
    {\itshape\grothFoldertitle}\\[4pt]
    {\small feuillets \grothFirst--\grothLast\ (lot \grothBatch)}\\[2pt]
    {\footnotesize\color{gray!60!black}Transcription établie d'après le fac-similé numérisé
     par l'Université de Montpellier.\\
     Les lectures incertaines sont soulignées, les illisibles notées [\,\ldots\,],
     les ajouts entre crochets.}
  \end{center}
  \vspace{1em}}

\endinput


\folder{161-3}
\batch{3}
\pages{41}{54}
\dating{[vers 1963-1973]}
\watermark{Édition de démonstration}
\foldertitle{Topos : notes manuscrites (s.d.)}

\begin{document}

\page{41}
\ldots\ par \ill{} \ill{} des corps). La question a une solution affirmative ssi
tout \struck{anneau} $\mathcal{C}$ ($\in \mathrm{Ob}(\mathrm{Ann})$) tel que, pour
tout $(u : A \to B) \in M$, $\operatorname{Hom}(B, \mathcal{C}) \xrightarrow{\ \sim\ }
\operatorname{Hom}(A, \mathcal{C})$ bij., $\mathcal{C}$ est un \emph{corps}.
\note{le début de la phrase est sur la page précédente, qui n'appartient pas à ce
lot}

Prenons donc d'abord \uncertain{$\lambda = \mathcal{E}\lambda_{0}$}, dans les
\ill{} des $S^{T}_{\lambda}$ (« $R^{T}_{\lambda}$ »), et les anneaux de polynômes
sur $\mathbb{Z}$ (i.e.\ une infinité de variables) $\bigl[$ \uncertain{vgr.} les
\struck{équations} relations (\ill{} aux solutions) \struck{affines} \ill{} types
en types finis $E^{\mathbb{Z}}_{\mathbb{Z}}$ sur $\mathbb{Z}$ $\bigr]$.

Si $u : X \to Y$ est dans $M^{\mathcal{C}}$, on \uncertain{sait} que
$m = u$, et que c'est une \ill{} \uncertain{fibrée} par fibre (Alors théorème de
Zariski) donc que c'est une \ill{} \struck{\ill{}} \uncertain{bien fermée}. Donc
$M'$ est réduit aux isomorphismes, donc $M$ aussi ; on aurait donc que
$(\mathrm{Ann})^{*} = \mathrm{Ann}$, absurde !
\note{$X$ et $Y$ portent chacun en dessous la mention $E^{m}$, et $X$ en outre un
$\lambda$}

b) Supposons \struck{\ill{}} que \uncertain{$\lambda = (\underleftarrow{\Sigma}, \underrightarrow{\lambda} = \emptyset)$}, et que $S^{T}_{\lambda}$ contient, en
plus \struck{\ill{}} de $\mathbb{Z}[t]$, l'anneau $\mathbb{Z}[t,t^{-1}]$ (ce qui
est vrai si $\underleftarrow{\lambda}$ = ens.\ de tous les diagrammes finis
\ill{} un seul cas $S^{T}_{\lambda}$ = ens.\ des anneaux de types finis sur
$\mathbb{Z}$). Alors la flèche canonique
\[
u_{0} : \mathbb{Z}[t] \longrightarrow \mathbb{Z}[t,t^{-1}] \times \mathbb{Z}
\qquad \Bigl( \simeq \mathbb{Z}[x,y] \big/
\bigl(x(xy-1),\, y(xy-1)\bigr) \Bigr)
\]
\note{le quotient est d'abord écrit avec un autre système de générateurs, rayé}
qui applique $t$ dans $(t,0)$, est dans $M$, et plus précisément, pour tout
anneau $\mathcal{C}$, le transformé de $u_{0}$ par $\operatorname{Hom}(-,
\mathcal{C})$ est l'application (\uncertain{surjective}) \struck{\ill{}} \ill{}
\[
\mathcal{C}^{*} \times \{e\} \longrightarrow \mathcal{C}
\]
de $\mathcal{C}$ en $\mathcal{C}_{1} \times \mathcal{C}_{2}$ (i.e.\ si
$\mathcal{X} = \operatorname{Spec} \mathcal{C}$, \ill{} $2 \simeq 2_{1} \amalg
2_{2}$) et : à \uncertain{un élément} inversible $\varepsilon$ de
$\mathcal{C}_{1}^{\#}$, associe l'élément $(\varepsilon, 0)$ de $\mathcal{C}$.
\note{l'annotation au-dessus de la ligne se lit « on aurait dans
$\widetilde{\mathcal{X}}$ »}

Cette application est bijective ssi $\mathcal{C}$ est un corps
$\bigl[$ puisque d'abord que \struck{$\mathbb{Z}/$} $\operatorname{Spec}
\mathcal{C}$ est connexe i.e.\ $\mathcal{C}$ n'est pas un produit $\mathcal{C}_{1}
\times \mathcal{C}_{2}$ avec $\mathcal{C}_{1}, \mathcal{C}_{2} \neq 0$
\note{la page s'arrête ici, le crochet reste ouvert}

\page{42}
c) Soit $G$ un groupe d'un topos, topos classifiant $\underline{B_{G}}$ (il
« classifie » les torseurs sous $G$ --- mais \ill{})

d) \emph{Topos induit} $E/X$ (on identifie \ill{} $E$ avec son objet
\ill{} ; $E_{/X}$ ou $X$ \ldots)

\textbf{3) Morphismes de topos : \emph{forment une catégorie}.}

\begin{tikzcd}
\underline{\mathrm{Mor}\,\mathrm{Top}}(X,Y) \arrow[r, hook, "\text{pl.\ fid.}"']
& \underline{\operatorname{Hom}}\,\underline{\mathrm{cat}}(X,Y)
\end{tikzcd}
\[
f \longmapsto f_{*}
\]

\marginal{\ill{} \ill{} des flèches}

La composition des morphismes est un \ill{} \struck{morphismes} de catégories,
m.\ associées. Les topos forment une « 2-catégorie » (\ill{} une
\uncertain{2-catégorie}). Notion d'équivalence de topos : $f : X \to Y$ est une
équivalence ssi $\exists\, g : Y \to X$ t.q.\ $gf \simeq \mathrm{id}_{X}$,
$fg \simeq \mathrm{id}_{Y}$.

Ici, c'est la notion habituelle d'équivalence de catégories. (C'est la notion qui
remplace la notion d'isomorphisme ($=$ homéomorphisme) pour les espaces
topologiques.

\marginal{Attention : pour l'instant, \ill{} d'isom\ill{} pour \ill{} de topos}

\textbf{4) Exemples de morphismes de topos}

a) Cas de $\mathrm{Top}(X)$ : déjà étudié : $\mathrm{Top}(X)$ dépend de $X$ de
façon covariante. On les identifie.

($\mathrm{Top}$ ord) \struck{\ill{}} peut être transformé en 2-catégorie. C'est un
foncteur 2-fidèle de \ill{} dans la 2-catégorie $(\mathrm{Top})$.

d$'$) (Cas de $\mathrm{Top}(X,G)$) (il n'y a plus de \uncertain{pleine fidélité})

b) $\bigl[$ Morphismes de $E$ \struck{dans} $\mathrm{Top}(X)$ : pareil que
\struck{foutes} \ill{} applications
\[
\mathrm{Ouv}(X) \longrightarrow \text{sous-objets de } e_{E}
\quad (\text{ouverts de } E),
\]
commutant aux $\sup$ quelconques et aux $\cap$ finies $\bigr]$
\note{sous la ligne : « ($\sim$ tous) $D$ »}

\emph{Ex.} Si $X$ est un \emph{doublet} $\bullet\!\!-\!\!\bullet$
($X = \{x,y\}$, \struck{$x \in \overline{\{y\}}$, $\notin \{x\}$, $\{x\}$}
$x$ fermé, $x \neq y$, $x \in \overline{\{y\}}$), alors le domaine de $E$ dans
$\mathrm{Top}(X)$ \ill{} : la donnée d'un ouvert de $E$ (\struck{le topos des
doublets} est topos classifiant pour l'objet de structure : sous-objet de l'objet
final). \emph{NB} Tout espace top.\ se \ill{} plonger dans un produit de
doublets $D^{I}$ (i.e.\ un \ill{} produit $J = \mathrm{Ouv}(X)$).

\page{43}
\begin{tikzcd}[column sep=small, row sep=small]
& S' \arrow[dr, dashed] & \\
S \arrow[ur] \arrow[rr] & & R
\end{tikzcd}

\begin{tikzcd}[column sep=large, nodes={font=\scriptsize}]
\operatorname{Hom}_{\underrightarrow{\lambda}}(R, \mathbb{Z}) \arrow[r, "\text{fid.}"]
& \operatorname{Hom}_{\underrightarrow{\lambda}}(S', \mathbb{Z}) \arrow[r, hook, "\text{pl.\ fid.}"]
& \operatorname{Hom}_{\underrightarrow{\lambda}}(S, \mathbb{Z})
\end{tikzcd}

\marginal{$T(\mathrm{Ens})$ ; $R^{\lambda}_{T}$ en fonction de $T(\mathrm{Ens})$ ?}

$\lambda \subset \lambda'$

\begin{tikzcd}[column sep=small, row sep=small, nodes={font=\scriptsize}]
(\mathrm{Cat}_{\lambda'}) \arrow[r] & (\mathrm{Cat}_{\lambda}) \\
\mathrm{Cat} \supset \mathrm{cofib}(\mathrm{Cat}_{\lambda}) \arrow[r] \arrow[d]
& \mathrm{Cat} \supset \mathrm{cofib}(\mathrm{Cat}_{\lambda'}) \\
(\lambda\text{-types}) \arrow[r, dashed, "?"] \arrow[d]
& (\lambda'\text{-types}) \arrow[u] \\
(\mathrm{Cat}_{\lambda}) \arrow[r, dashed]
& (\mathrm{Cat}_{\lambda'}) \arrow[u]
\end{tikzcd}
\note{la lecture de $\mathrm{cofib}$ est douteuse aux deux sommets ; la flèche
horizontale du bas est tracée deux fois, la première rayée, la seconde en
pointillé}

\marginal{$\lambda'$-catégorie \struck{engendrée} \ill{} d'une $\lambda$-catégorie}

\begin{tikzcd}[row sep=small]
C \arrow[r, hook] & (\check{C})^{\circ}
\end{tikzcd}
\note{sous ce plongement la page porte encore $\bigvee$ puis $X_{i}$, dont le rôle
n'est pas explicité}

\[
\varinjlim\nolimits_{\lambda}(R, C^{\circ})^{\circ} \simeq
\varinjlim\nolimits_{\lambda}(R^{\circ}, C)
\]

\struck{\ill{}} $T(C)$
\begin{tikzcd}[column sep=small, row sep=small, nodes={font=\scriptsize}]
\operatorname{Hom}_{\lambda}(R, C) \arrow[r, hook, "\text{2-cart}"] \arrow[d, "\varphi_{C}"]
& \operatorname{Hom}_{\lambda}(R, \widehat{C}\,) \arrow[d, "\varphi_{\widehat{C}}"] \\
C^{I} \arrow[r, hook] & \widehat{C}^{\,I}
\end{tikzcd}
\note{au-dessus du carré la page porte $C^{\circ\,I}$ et $C \to C^{I}$}

À droite, pour $C \hookrightarrow C'$ ($\lambda$-fidèle, pl.\ fid.) :
\begin{tikzcd}[column sep=small, row sep=small]
T_{\lambda}(C) \arrow[r, hook] \arrow[d] & T_{\lambda}(C') \arrow[d] \\
C^{+} \arrow[r, hook] & C'^{+}
\end{tikzcd}

\[
T(\mathrm{Ens}) \simeq \operatorname{Hom}_{\lambda}(R, (\mathrm{Ens})), \qquad
R_{T} \Longrightarrow \operatorname{Hom}(T(\mathrm{Ens}), \mathrm{Ens})
\]

\begin{itemize}
\item 1) \uncertain{Rigidité} : $\widehat{C}$   $R_{T}
\xrightarrow{\ \sim\ } \operatorname{Hom}_{\lambda}\bigl((T(\mathrm{Ens}))_{\circ},
(\mathrm{Ens})\bigr)$ (quelque \ill{})
\item 2) Existence des $\varinjlim$ \struck{\ill{}} de type $\lambda$ dans $T(C)$
(\ill{} cela $\underrightarrow{\lambda}$ et $\underleftarrow{\lambda}$)
\item 3) Existence des $\varprojlim$ dans les $T(C)$ ni \ldots\ (?)
\item 4) Adjoint de $\varphi : T(\mathrm{Ens}) \to (\mathrm{Ens})^{I}$ dans
$T(\mathrm{Ens})$ si $\lambda$ est formé de diag.\ \emph{finis}
\item 5) $R_{T} \simeq \operatorname{Hom}_{\lambda}\bigl((T(\mathrm{Ens}))_{\circ},
(\mathrm{Ens})\bigr)$ ? ?   ($\hookrightarrow$ objets de t.f.)
\end{itemize}

\page{44}
\emph{NB} Si l'ens.\ des parties de $E$ est \struck{\ill{}} isomorphe (pour
la structure ordonnée) à l'ens.\ des ouverts $\mathrm{Ouv}(X)$ pour quelque $X$,
alors $E \to$ \struck{\ill{}} $[$i.e.\ $E \to \mathrm{Top}(X)]$, et c'est
universel pour les morphismes de $E$ dans un espace topologique.

\struck{c) $\mathrm{Top}(X,G)$ pour \ill{} les morphismes de \ill{}
$E \to \operatorname{Hom}(\mathcal{A}(X), \mathrm{Ens})$ \ill{}
$\to \operatorname{Hom}(E, (\mathrm{Ens}))$}

c) Morphismes d'un pt dans \ill{} topos $E$ : catégorie
\struck{$\underline{\mathrm{Pt}}(E)$} $\underline{\mathrm{Pt}}(E)^{\circ}$

d) \emph{Catégories} $\widehat{C}$ : soit \struck{\ill{}} $u : C \to
\widehat{C}'$ définit un foncteur en sens inverse $F' \mapsto F \circ u$

\begin{tikzcd}
\widehat{C}' \arrow[r, "u^{*}"] & \widehat{C}
\end{tikzcd}

\struck{qui a des adjoints}, notés $u_{!}$ et $u^{!}$ qui commutent aux
$\varprojlim$ quelc.\ et \ill{} $\varinjlim$ quelc., donc est \ill{}
\struck{foncteur $u^{*}$}. Les foncteurs image inverse pour un $u : \widehat{C}
\to \widehat{C}'$ sont $u_{!}$, $u^{*}$, $u_{*}$.

\marginal{dans $\widehat{C}$ \ill{} on retrouve \ill{} $C$ \ill{} covariant}

Elle peut avoir l'adjoint à droite $u_{*}$ ; il a un adjoint $\varphi$, qu'on note
$u_{!}$ \ldots

\emph{Ex.} $B_{G}$ avec $G$ variable, qui est aussi exemple de 4.1.

e) Morphismes \ill{} de topos $\widehat{C} \to E$ :
\[
E \xrightarrow[\ \varprojlim \text{finies}\ ]{\ \varprojlim\ }
\operatorname{Hom}(C^{\circ}, \mathrm{Ens})
\qquad \Longleftrightarrow \qquad
C^{\circ} \longrightarrow \operatorname{Hom}_{\varprojlim\,\text{finies}}
(E, (\mathrm{Ens})) = \underline{\mathrm{Pt}}(E)^{\circ}
\]
\[
\operatorname{Hom}_{\mathrm{top}}(\widehat{C}, E) \simeq
\operatorname{Hom}(C, \underline{\mathrm{Pt}}(E))
\]
\note{le second membre est d'abord écrit
\struck{$\operatorname{Hom}(C^{\circ}, \underline{\mathrm{Pt}}(E)^{\circ})$}}
si $E = \widehat{C}'$, on trouve
\[
\operatorname{Hom}_{\mathrm{top}}(\widehat{C}, \widehat{C}') \simeq
\operatorname{Hom}(C, \mathrm{Pro}\,C')
\qquad \bigl(\operatorname{Hom}(C, C')\ \text{pl.\ fid.}\bigr)
\]
\note{$\mathrm{Pro}$ est précédé d'un \struck{$\mathrm{Pt}$} rayé}
si $C$ stable par $\varprojlim$ finies, on obtient
\[
\operatorname{Hom}_{\mathrm{top}}(E, \widehat{C}) \simeq
\operatorname{Hom}_{\mathrm{ex.\,g.}}(C, E)^{\circ}
\]

\page{45}
$\lambda = (\underleftarrow{\lambda}, \emptyset)$

\marginal{muni des foncteurs de cat.\ \uncertain{2}-abstraites
$b^{T}_{C} : T(C) \to C^{I}$}

\textbf{1)} $T$ catégorie $\lambda$-abstraite sur $(\mathrm{Cat}_{\lambda})$, on
suppose que \struck{\ill{} $C$}

a) Pour tout foncteur pl.\ fidèle \struck{si} $C \to C'$ dans
$\mathrm{Cat}_{\lambda}$, $T(i) : T(C) \to T(C')$ pl.\ fid.

b) et de même
\begin{tikzcd}[row sep=small]
T(C) \arrow[r] \arrow[d, "b^{T}_{C}"] & T(C') \arrow[d, "b^{T}_{C'}"] \\
C^{I} \arrow[r] & C'^{I}
\end{tikzcd}
est 2-cartésien $[$il \ill{} suffit pour $C \hookrightarrow \widehat{C}\,]$

\marginal{Attention : \ill{} $\widehat{C}$ n'est pas une $\mathcal{U}$-catégorie !}

b) Pour tout $C \in \mathrm{Ob}\,\mathrm{Cat}_{\lambda}$, le foncteur canonique
\[
T(\widehat{C}\,) \longrightarrow \operatorname{Hom}_{\lambda}(C^{\circ}, S)
\qquad \bigl( S \overset{\text{déf}}{=} T(\mathrm{Ens}) \bigr)
\]
est une équivalence de catégories.

Alors $T(C)$ se reconstitue à partir de la donnée de $S$ et de
\[
\beta = (b_{\mathrm{Ens}}) : S \longrightarrow (\mathrm{Ens})^{I},
\qquad \bigl( \beta = (\beta_{i})_{i \in I},\ \beta_{i} : S \to (\mathrm{Ens})
\bigr)
\]
$T(C)$ = sous-catégorie pleine de $\operatorname{Hom}_{\lambda}(C^{\circ}, S)$
formée des $\varphi : C^{\circ} \to S$ tels que $\forall\, i \in I$,
$\beta_{i} \circ \varphi$ soit représentable.

\marginal{Mais si $f : C \to C'$ est un $\lambda$-foncteur, comment définir
$T(f) : T(C) \to T(C')$ en fonction de $S$, $(\beta_{i})$ ?}

\struck{\ill{}} La sous-catégorie pleine $R$ de $\operatorname{Hom}_{\lambda}(S,
(\mathrm{Ens}))$ formée des foncteurs $\psi$ tels que pour tout $C \in
\mathrm{Ob}\,\mathrm{Cat}_{\lambda}$ et $\xi \in T(C)$, le composé
\[
\psi \circ \xi : C^{\circ} \longrightarrow S \xrightarrow{\ \psi\ }
(\mathrm{Ens})
\]
soit représentable, est stable par $\varprojlim$ de type $\in
\underleftarrow{\lambda}$, \struck{À quelle condition} et pour tout
$C \in \mathrm{Ob}\,\mathrm{Cat}_{\lambda}$, on a un accouplement
\[
T(C) \times R \longrightarrow C
\]
commutant aux $\varprojlim$ de type $\underleftarrow{\lambda}$ sur $R$, d'où
\[
T(C) \longrightarrow \operatorname{Hom}_{\lambda}(R, C)
\]
(fonct.\ de cat.\ 2-\ill{} sur $(\lambda)$)

À quelle condition cela est-il une équivalence de catégories ? Il est nécessaire
que

a) $R$ soit \struck{formé de foncteurs} repr.\ i.e.\ $R \simeq \Sigma^{\circ}$, où
$\Sigma$ est une sous-catégorie pleine de $S$ $[$dans laquelle les
$\varinjlim$ (\uncertain{inductibles}) de type $\lambda$ existent dans $S$, et
\ill{} : $\Sigma\,]$. Donc la donnée de $R$ équivaut à la donnée de
$\Sigma \subset S$.

b) $\Sigma$ engendre $S$ $[\, S \to \operatorname{Hom}(\Sigma^{\circ},
(\mathrm{Ens})) = \widehat{\Sigma}$ pl.\ fidèle $]$

\page{46}
e) \emph{Morphismes dans topos ponctuel} $\mathrm{Top}(1)$ : objet
constant\add{s}, foncteur sections.

f) Morphismes du topos vide dans un topos. Morphisme dans le topos vide
$E \to$ \struck{$\mathrm{Ens}$} $\mathrm{Top}(\emptyset)$ : c'est ssi
$E \simeq \mathrm{Top}(\emptyset)$, i.e.\ $E$ équivalent à cat.\ ponctuelle (on
aurait tous les objets isomorphes, on aurait \ill{} morphisme
$\emptyset_{E} \to e_{E}$ un \ill{} \ldots)

g) Classifiant $B_{G}$ ($G$ groupe d'un topos) \ill{} variable \ill{} objet $G$.
\[
E \rightleftarrows B_{G} \quad \text{correspond à} \quad e \rightleftarrows S
\]

h) Morphismes $E_{/X} \longrightarrow E$ d'\emph{induction}
\[
\begin{cases}
u^{*}(Y) = Y_{X} & \\
u_{!}(Y) = Y & (\text{oubli de plus structure}) \\
u_{*}(Y) & \text{plus compliqué}
\end{cases}
\]
Ici \struck{\ill{}} $u_{!}$, $u^{*}$, $u_{*}$ ; $u^{*}$ commute aux
$\varprojlim$ quelconques \ldots

\marginal{morphismes $E' \to E$, $2^{\circ}$ des \ill{} \ill{} locaux \ill{} ($=$
\ill{})}

\marginal{foncteurs de localisation, morphismes du localisé}

$E_{/X}$ \struck{\ill{}} comme foncteur en $X$ : on \ill{} que
$\operatorname{Hom}_{\mathrm{top}_{E}}(E_{/X}, E_{/Y}) \simeq$ catégorie discrète
de « $u_{!}$ » $\operatorname{Hom}(X,Y)$

i) \emph{Topos somme de topos}
\note{le titre seul ; le reste de la page est vierge}

\page{47}
Donc $\beta$ $(\beta_{i})$ s'interprète en termes de cat.\ des \ill{},
(\uncertain{vgr.} constr.\ \ill{})

$\lambda$ = ens.\ des types de diagrammes \emph{finis}

\begin{itemize}
\item $\Sigma$ catégorie stable par $\varprojlim$ de type $\lambda$
\item $R = \Sigma^{\circ}$
\item $S = \operatorname{Hom}_{\underrightarrow{\lambda}}(R, \mathrm{Ens})
\supset \widehat{\Sigma} \supset \Sigma$, avec un \struck{$\mathrm{Ens}$}
rayé après la parenthèse
\end{itemize}

Je dis que

a) $S$ admet des $\varinjlim$ quelc.

b) Pour les catégories $\mathcal{T}$ ayant des $\varinjlim$ quelc., on a
\[
\operatorname{Hom}_{\underrightarrow{\lambda}}(\Sigma, \mathcal{T})
\xrightarrow{\ \sim\ } \operatorname{Hom}_{\lambda'}(S, \mathcal{T})
\]
où $\operatorname{Hom}_{\lambda'}$ désigne les foncteurs commutant aux
$\varinjlim$ quelconques ; et
\[
\operatorname{Hom}_{\underleftarrow{\lambda}}(R, \mathcal{T}^{\circ})
\qquad \parallel
\]
\ill{} $(\mathcal{T}^{\circ}) \xrightarrow{\ \sim\ }
\operatorname{Hom}_{\lambda'}(S, \mathcal{T})$

$S$ engendré par $\varinjlim$ par $\Sigma$.

\page{48}
\note{les pages 48 puis 52 à 54 sont rédigées en anglais ; la transcription suit
la langue de la page}

Comparison between topoi $\widehat{C}$ and $\mathrm{Top}(X)$.

$\widehat{C}$ is $\simeq$ to a $\mathrm{Top}(X)$ iff $C$ is an ordered category
(i.e.\ $\forall\, x \in \mathrm{Ob}\,C$, $x \to e_{\widehat{C}}$ is a
monomorphism). When this is so, we can take $X$ to be $C$ with the topology
generated by the sets $O_{x} = \{ y \in C \mid y \leq x \}$, i.e.\ arbitrary
unions (of intersections) of such sets: in fact \struck{\ill{}} the $O_{x}$'s form
a \emph{basis} of the topology $\mathcal{T}$, and open \struck{On the other
hand, if $X$ is an} sets \struck{are those which} the \struck{\ill{}} sieve's
\[
\bigl( \forall\, x \in O,\ (y \leq x \text{ and }
x \in O \implies y \in O) \bigr).
\]
\note{un \struck{$\Longleftrightarrow$} rayé précède la seconde parenthèse}
\note{sous cette ligne : « and $O = \bigcup O_{x}$ »}

If $F$ is a \struck{functor on $C$, then} $F(O) \Longrightarrow \prod F(O_{x})
\rightrightarrows \prod F$. Thus sheaves on $C$ are just the sheaves \ill{} this
basis of open sets, i.e.\ contravariant functors $O_{x} \mapsto F(O_{x}) = F(x)$
from $C'$ \struck{\ill{}} (category of the $O_{x}$, $O_{y} \hookrightarrow O_{x}$
iff $y \leq x$) $\simeq C$ to sets, satisfying the sheaf property
\struck{\ill{}} with respect to \struck{\ill{}} coverings of $O_{j}$, but these
are trivial $[$for if $O_{x_{i}}$ cover $O_{x}$, $\exists\, i$ s.th.\
$x \in O_{x_{i}}$, i.e.\ \struck{$x_{i}$} $x \leq x_{i}$, hence $x \approx x_{i}$
(as $x_{i} \leq x$)$]$ then we'll get arbitrary products.

Conversely, if $X$ is a \uncertain{sober} space, then $\mathrm{Top}(X)$ is
$\simeq$ ($\widehat{C}$ for some \ill{}) iff the \struck{\ill{}} essential points
of $X$ (i.e.\ those which have a smallest open neighbourhood --- or those for
which the fiber functor \struck{commutes with products}) are very dense $[$and we
may take for $C$ the category of these pts$)$
\note{la parenthèse fermante est écrite à la place du crochet}

\page{49}
\begin{tikzcd}[column sep=huge, row sep=small]
C \arrow[r, "f"] & C' \\
C & C' \arrow[l, dashed, "g"'] \\
C \arrow[r, dashed, "h"] & C'
\end{tikzcd}
\note{sur la page les trois flèches sont superposées entre $C$ et $C'$, $f$
pleine, $g$ et $h$ en pointillé}

\begin{tikzcd}[column sep=huge, row sep=small]
\widehat{C} \arrow[r, "f_{!}"] & \widehat{C}' \\
\widehat{C} & \widehat{C}' \arrow[l, "f^{*}"'] \\
\widehat{C} \arrow[r, "f_{*}"] & \widehat{C}'
\end{tikzcd}
\note{sur la page les trois foncteurs sont superposés entre les deux mêmes
sommets}

$f_{!}$ prolonge $f$ ; $f^{*}$ prolonge $g$ (si $g$ existe) ; $f_{*}$ prolonge $h$
(si $h$ existe).

Si $f$ commute aux $\varprojlim$ finies, alors $f_{!}$ aussi ;
\struck{\ill{}} $\varphi^{*}$, $\varphi : \widehat{C}' \to \widehat{C}$ un
morphisme de topos.
\[
\widehat{f}_{!} = \varphi^{*}, \qquad
\widehat{f}^{\,*} = \varphi_{*}, \qquad
\widehat{f}_{*} = \varphi^{!}
\]

$\widehat{f} : \widehat{C} \to \widehat{C}'$ morphisme de topos, avec
$\widehat{f}^{\,*} = f^{*}$, $\widehat{f}_{*} = f_{*}$, $\widehat{f}_{!} = f_{!}$.

\[
\operatorname{Hom}_{\widehat{C}'}(f_{!}(F), F') \simeq
\operatorname{Hom}_{\widehat{C}}(F, f^{*}F')
\]
si $F = X \in \mathrm{Ob}\,C$, cela signifie
\[
\operatorname{Hom}_{\widehat{C}'}(f_{!}(X), F') = f^{*}(F')(X) = F'(f(x))
= \operatorname{Hom}_{\widehat{C}'}(f(x), F')
\]
i.e.\ $f_{!}(x) = f(x)$.
\[
\operatorname{Hom}_{\widehat{C}'}(F', f_{*}F) \simeq
\operatorname{Hom}_{\widehat{C}}(f^{*}F', F)
\]
si $F = x \in \mathrm{Ob}\,C$, cela signifie
\[
\operatorname{Hom}_{\widehat{C}'}(F', f_{*}(x)) =
\operatorname{Hom}_{\widehat{C}}(f^{*}F', x)
\]
et faisant $F' = x' \in \mathrm{Ob}\,C'$ :
\struck{\ill{}}
\[
f_{*}(x)(x') = \operatorname{Hom}_{\widehat{C}}
(x' \circ f^{\circ}, x), \qquad
f_{*}(x) = \bigl( x' \longmapsto \operatorname{Hom}_{\widehat{C}}
(x' \circ f^{\circ}, x) \bigr)
\]
\[
x' \circ f^{\circ} = \bigl( y \longmapsto \operatorname{Hom}(f(y), x') \bigr)
\ \parallel\ \operatorname{Hom}(y, g\,x')
\]
\[
f_{*}(x) = \bigl( x' \longmapsto \operatorname{Hom}_{C}(g(x'), x) \bigr)
\ \wr\ \operatorname{Hom}_{C}(x', h(x))
\]

La page porte encore, à part, $(\mathbb{Z}, G)$ et $B_{G} \to B_{e}$.
\note{ce fragment annonce les pp.~50-51}

\page{50}
\section*{Exemples de morphismes et de foncteurs de localisation}

$E$, $G$ ; $B_{G}$ ; $H \subset G$ ; $X = G/H$ ; $G \in \mathrm{Ob}\,B_{G}$, muni
de $\varepsilon \in \Gamma(X)$.

\begin{tikzcd}[row sep=large]
(B_{G})_{/X} \arrow[r, "\varphi\ \approx"] \arrow[d, "j_{X}"]
& B_{H} \arrow[d, "B_{i}"] \\
B_{G} & B_{G}
\end{tikzcd}

\begin{tikzcd}[row sep=small]
E \arrow[r] & X \\
F \arrow[u] \arrow[r, dashed] & e_{E} \arrow[u, "\varepsilon"']
\end{tikzcd}

$\varphi(E) = E \times_{X} (e_{E}, \varepsilon)$,
$\psi(F) = G \wedge_{H} F$

\[
\varphi(j^{*}(E)) = (E \times X) \times_{X} (e_{E}, \varepsilon) \simeq E
\]

\marginal{oubli partiel des opérations de $G$ --- ou plutôt de $H$}

Donc $j_{X}$, via l'équivalence $\varphi$, s'interprète comme $B_{i}$. On a
$j_{X!}$ s'interprète comme $F \mapsto G \wedge_{H} F$, et $j_{X*}$ s'interprète
comme $F \mapsto \operatorname{Hom}_{H}(G, F)$.

\marginal{plus généralement, si $Z \in \mathrm{Ob}\,B_{G}$, \ill{} foncteur
« oubli » $(B_{G})_{/(Z,G)} \to B_{G} \to B_{H}$ \ill{} $(B_{H})_{(Z,H)} \to B_{H}$
\ill{} sont des foncteurs de localisation}

Cas des groupes discrets ($E = (\mathrm{Ens})$) : si $G$ est le groupe fondamental
d'un espace topologique $Z$ connexe, loc.\ connexe et loc.\ 1-connexe
\struck{et \ill{} plus connexe}, $B_{G} \simeq$ revêtements de $Z$, et le domaine
de $X$ s'identifie à la donnée d'un revêtement \struck{\ill{}} de $Z$, pointé au
dessus de $e$ (i.e.\ un $x \in X$).

\marginal{théorie de Galois pour \uncertain{$Z$} locaux}

Alors la \struck{\ill{}} équivalence $\varphi$ \struck{\ill{}} s'interprète comme
l'équivalence entre $\mathrm{Rev}(X)$ (pour $(B_{G})_{/X}$) et
\struck{la théorie de} $B_{H}$, dont c'est la théorie de Galois pour $X$ et le pt
base $x$ (compte tenu que $\pi_{1}(X, x) \simeq H \subset G$).

TSVP

\page{51}
Soit de \struck{même} \struck{$\mathbb{Z}$} $X$ un espace topologique avec
\struck{\ill{}} un groupe \add{discret $G$} qui opère, d'où
\struck{$\mathrm{Top}$} $\mathrm{Top}(\mathbb{Z}, G)$ ; si $H$ est un sous-groupe
de $G$, on a un foncteur « restriction des opérations à $H$ »
\[
\mathrm{Top}(\mathbb{Z}, G) \longrightarrow \mathrm{Top}(\mathbb{Z}, H)
\]
qui s'identifie à un foncteur de localisation, à savoir : l'objet \ill{} pour
$X = G/H$, on prend $\mathbb{Z}' = \mathbb{Z} \times X$ avec opérations
diagonales de $G$, et on considère le relatif $\mathbb{Z}'$ sur $\mathbb{Z}$
(c.-à-d.\ $\mathrm{Top}(\mathbb{Z})$) défini par l'élément marqué $x$ de $X$. Cela
dit, on \struck{prend} le foncteur
\[
\mathrm{Top}(\mathbb{Z}, G)_{/\mathbb{Z}'} \longrightarrow
\mathrm{Top}(\mathbb{Z}, H), \qquad E \longmapsto E \times_{\mathbb{Z}'}(X, e)
\]
est une équivalence de catégories \ldots\ (En fait, $\mathbb{Z}'$ est un
« revêtement topologique » de $\mathbb{Z}$ \ldots)
\note{la phrase est laissée bancale sur la page : « on le foncteur \ldots\ est une
équivalence »}

\page{52}
\section*{Sites}

\begin{itemize}
\item (1) Notion of localisation: yoga.
\item (2) localisation \add{data} (in a category (a
« \emph{topology} »))
\end{itemize}

a) In terms of « coverings » (when fibered products exist)
\begin{itemize}
\item Cov 1   base change
\item Cov 2   composition
\item Cov 3   identity
\item $[$ Cov 4   saturation $]$
\end{itemize}
\note{une accolade réunit Cov 1 à Cov 3 : « postulating, assuming only the
$X_{\alpha} \to X$ squarable »}

b) In terms of « sieves » \add{covering}
\begin{itemize}
\item T1   base change
\item T2   local nature
\item T3   \struck{\ill{}}
\end{itemize}

\begin{itemize}
\item order relation ($=$ \struck{\ill{}} inclusion)
\item Intersection of topologies / \struck{\ill{}} supremum
\item Topology generated by sieves, or families of morphisms
\end{itemize}

\marginal{notion of a sieve of $C$ (two versions) of $C_{/X}$ (a $X$).
Operations: base change, intersection, order relation, intersection, union; sieve
generated by a covering, and interpretation of operations \struck{\ill{}} in terms
of these}

\[
\operatorname{Hom}_{\widehat{C}}(R, F) = \varprojlim\nolimits_{C/R} F(X)
\simeq \operatorname{Ker}\Bigl( F(X_{\alpha}) \rightrightarrows
\prod\nolimits_{\alpha, \beta} F(X_{\alpha} \times_{X} X_{\beta}) \Bigr)
\]
\note{sous la formule : « family $J_{\rho}(X)$ »}

If Cov 1 to Cov 3 is satisfied, then the sieves generated by the objects of
$\mathrm{Cov}(X)$, are cofinal in $J_{\mathcal{F}}(X)$ ($\mathcal{T}$ top.\
$J$-generated).

(3) \emph{Sheaves} (and separated \add{pre}sheaves)
\[
\operatorname{Hom}(X, F) \longrightarrow \operatorname{Hom}(R, F)
\quad \text{mono. (biun.)} \quad \text{for } \forall\, R \subset X,\
R \in J(X)
\]
\[
\Bigl[ \operatorname{Hom}(X, F) \longrightarrow \operatorname{Ker}
\Bigl( \prod F(X_{\alpha}) \rightrightarrows F(X_{\alpha} \times_{X} X_{\beta})
\Bigr) \Bigr] \quad \text{mono. / biun.} \quad
\forall\, (X_{\alpha} \to X) \in \mathrm{Cov}(X)
\]

Topology defined by a family $\Phi$ of presheaves, to become sheaves
(resp.\ separated presheaves): topology $T_{\Phi}$ where
$(R \hookrightarrow X) \in T_{\Phi}(X)$ iff $\forall\, X' \to X$ in $C$, if
$R' = R \times_{X} X'$, $\forall\, F \in \Phi$, $F(X') \to F(R')$ bij.\ (inj.).
This topology is the finest \struck{\ill{}} among those turning $F \in \Phi$ into
sheaves (resp.\ presheaves).

\emph{Cor} If T1 is defined by a family of sieves, then cat.\ of sheaves
(resp.\ presheaves) for topology generated is easy $[$if T1 not satisfied one has
first to saturate by T1$]$

\page{53}
\note{la numérotation des théorèmes de cette page fait suite à celle de la
p.~54 : Th.~1 et Th.~2 y sont énoncés}

the objects of $\widetilde{C}$ \struck{are} sheaves (or only separated
presheaves). In fact, we get

\emph{Th.\ 3} (Giraud) $T \longrightarrow \widetilde{C}_{T}$ is a bijective
correspondence between topologies on $C$ and \add{strictly} full subcategories $E$
of $\widehat{C}$ s.th.\ the inclusion functor $i : E \to \widehat{C}$ has a left
adjoint which is left exact. It is order reversing ($T \leq T' \Longleftrightarrow
\widetilde{C}_{T'} \subset \widetilde{C}_{T}$).

\emph{Th.\ 5} Let $E$ be a category. Then $E \simeq \widetilde{C}$ for some
(small) $C$ iff $E$ satisfies a) b) c) d). We can simply choose
\struck{take for $C$ the class} \ill{} the canonical \struck{a category $C$ with a
topology} of a top.\ $T$ in $E$, (i.e.\ $C \subset \widehat{C}$). The choice
\struck{\ill{}} $\widetilde{C}_{T} \simeq E$ is equivalent to a
\struck{full subcategory of $E$} \struck{and the topology} \struck{\ill{} is
generating} fully faithful functor $C \to E$ such that the image of $C$ is
generating $E$; the topology of $C$ is \emph{induced} by the canonical
topology of $E$, via the given functor \ldots\ (for instance by
$C \xrightarrow{\ \varepsilon_{C}\ } E = \widehat{C}$)

\struck{Continuous functors between sites, induced topology} and:

\emph{Th.\ 4} \emph{Comparison lemma} $C \hookrightarrow C'$
($\mathcal{U}$-cat.) \ill{} de sites, $i'$ pleinement fidèle, top.\ de $C$
induite. Si tout objet de $C'$ peut être couvert par des objets de $C$, alors
$\widetilde{C}' \xrightarrow{\ \sim\ } \widetilde{C}$. (Il y a réciproque si $C$
est une $\mathcal{U}$-cat.\ dont la top.\ est moins fine que la top.\ can.)

\marginal{Analogue bien connu en top.\ ordinaire ; possibilité de définir des
$\widetilde{C}$ équivalents \ill{} de bien des façons différentes}

\page{54}
by taking \struck{as} $K'(X) =$ set of all inverse images of $R' \in K(X')$ by any
$X \to X'$ in $C$, then $F$ \struck{separated} sheaf (separated presheaf) for $T$
iff \struck{\ill{}} $\forall\, X \to X'$, $R' \in K(X')$, and taking
$R = X \times_{X'} R'$, $F(X) \to F(R)$ is inj.\ (bij.).

canonical top.\ $J_{T}$ and strict universal epimorphic families. Most usual
topologies are less fine than can., hence $C \to \widetilde{C}$ fully faithful.

\marginal{canonical topology is important in $E = \widetilde{C}$, where it will
mean that the covering families are those which have epimorphic \ill{}}

(4) \emph{Sheaf associated to a presheaf}
\[
\mathrm{id} \longrightarrow L, \qquad F \xrightarrow{\ \ell(F)\ } LF
\]
\[
LF(X) = \varinjlim\nolimits_{R \in T(X)} \operatorname{Hom}(R, F)
\]
\begin{itemize}
\item a) $L$ is left exact: $\widehat{C} \to \widehat{C}$
\item b) $\forall\, F$, $LF$ is separated
\item c) $F$ separated $\Longleftrightarrow F \to LF$ mono. If so, $LF$ is a
sheaf
\item d) $F \to LF$ is an isom.\ iff $F$ is a sheaf
\end{itemize}

\marginal{bounding the size of \ill{} $T(X)$ : \ill{} restrict to
$\mathcal{U}$-\ill{} \ill{} (\ill{} \ill{} or \ill{} de topoi \ill{} \ill{}
sites \ldots)}

\emph{Th.\ 1} Existence of $a : \widehat{C} \to \widetilde{C}$, left adjoint
of inclusion, and \struck{\ill{}} $a$ is left exact. Indeed,
$a = L \circ L$ \ldots

\emph{Cor} \ a) Projective limits exist in $\widetilde{C}$, and
$\widetilde{C} \subset \widehat{C}$ commutes with them.
b) Direct limits exist, $\varinjlim_{i}^{\widetilde{C}} F_{i} \simeq
a\bigl(\varinjlim_{i}^{\widehat{C}} F_{i}\bigr)$.
\note{d'où toutes les propriétés d'exactitude de (Ens) qui n'engagent que des
$\varprojlim$ finies et des $\varinjlim$ quelconques}

\emph{Cor} \ $F \simeq \varinjlim_{C/F} a(X)$ pour $F \in
\mathrm{Ob}\,\widetilde{C}$.

\emph{Th.\ 2} Consider $C \xrightarrow{\ \varepsilon_{C}\ } \widetilde{C}$
($X \mapsto a(X)$, left exact). Then

a) \struck{$a$ fully faithful}   $X_{\alpha} \to X$ is covering
$[$i.e.\ for every $F \in \mathrm{Ob}\,\widetilde{C}$, $F(X) \to \prod
F(X_{\alpha})$ is injective$]$ iff $\coprod \varepsilon_{C}(X_{\alpha})
\to \varepsilon_{C}(X)$ is epimorphic

b) $R \subset X$ is a covering sieve iff $\varepsilon_{C}(R)
\xrightarrow{\ \sim\ } \varepsilon_{C}(X)$.

Thus we recover topology $T$ from the knowledge of $\widetilde{C} \subset
\widehat{C}$; indeed, it is the \uncertain{finest} topology for which
\struck{\ill{}}
\note{la page s'arrête au milieu de la phrase}

\end{document}
